\chapter{总结与展望}

\section{研究工作总结}
本文的主要内容在前面章节中已经介绍，本章在此基础上对本文的研究工作进行总结:

(1) \textbf{方面级图文情感相关性语料构建与数据分析：}本研究从图文关系的视角出发，针对现有图文关系定义仅关注文本与图像的显式对齐或全局相关性，而忽略隐式情感关联，导致模型在复杂场景下易受干扰的问题，定义了基于情感表达的方面级跨模态相关性，并构建了一个新的方面级图文相关性数据集。在此基础上，本研究对该数据集进行了详细的统计分析，深入探究图文关系与情感之间的内在关联，为后续研究奠定了数据基础。同时，通过在多个数据集上将跨模态相关性信息融入多种多模态方面级情感分类模型进行对比实验，验证了所提出的方面级跨模态相关性的有效性。

(2) \textbf{基于图文情感线索引导的多模态方面级情感分析：}尽管前述工作验证了加入本文提出的方面级跨模态相关性能带来性能提升，但仅依靠粗粒度的相关性难以精确控制图像模态的贡献，且模型无法精确捕捉单模态中各自蕴含的情感线索，导致模型在情感分析场景下的鲁棒性不足。因此，本研究提出了基于图文情感线索引导的多模态方面级情感分析方法。针对现有方法对图文情感线索建模粒度不足的问题，该方法利用多模态大模型提取与方面相关的情感线索描述，将其作为补充信息编码至上下文；同时利用细粒度方面级跨模态相关性作为引导信号，控制图文情感线索与图文特征的细粒度融合，有效抑制无关或冲突线索的干扰，从而提升模型对图文情感的理解能力。

(3) \textbf{基于多角度图文相关性解耦的多模态方面级情感分析：}为了进一步提升模型对复杂图文的理解能力和情感识别能力，同时减少无关图像语义噪声的引入，本章提出了一种基于多角度图文相关性解耦的多模态方面级情感分析方法。为了实现对图像特征的精细化控制，本研究将方面级跨模态相关性进一步解耦为语义相关性与情感相关性两个独立维度进行建模。具体而言，模型分别提取图像中的语义特征和情感特征，并利用解耦后的多角度相关性概率作为门控引导信号；这一机制能够在有效过滤无关语义干扰的同时，针对性地增强关键情感特征的融合与利用。此外，为了夯实多角度相关性建模的基础，本研究引入了单模态情感监督辅助任务。该机制将情感线索从原本的输入端提示转化为训练阶段的监督信号，通过多任务学习促使模型主动习得鲁棒的单模态情感表征，从而显著提升模型在复杂场景下的情感分析能力与鲁棒性。


\section{下一步工作设想}

虽然本文在基于图文关系的多模态情感分析上取得一定研究成果，但后续工作还需要在以下方面进行进一步探索：

首先，为了进一步提升模型在处理反讽、隐喻等高阶语义复杂场景下的表现，后续研究将探索多模态大模型与思维链技术的深度融合。现有的解耦机制虽然有效过滤了噪声，但在解释“为何图文相关”的逻辑层面仍有提升空间。未来的工作计划利用大模型的生成能力，不仅仅提取情感线索，而是生成显式的推理路径，将隐式的跨模态情感关联转化为可解释的逻辑链条。通过知识蒸馏或参数高效微调技术，将这种深层的推理能力注入到本研究提出的轻量级模型中，从而在保持计算效率的同时，显著增强模型对于复杂图文语境的理解与推断能力。

其次，随着社交媒体内容形式的演变，将现有的静态图文分析框架扩展至视频多模态方面级情感分析是必然趋势。视频数据引入了时间维度与音频模态，带来了更高的数据冗余与对齐难度。考虑将方面级跨模态情感相关性的语义和情层面解耦的思想拓展至视频维度，设计一种动态的跨模态解耦机制。能够在从连续的视频帧序列与音频流中动态识别并分离出承载关键情感信息的情感片段，同时过滤掉仅包含背景语义的冗余帧与环境噪声。实现对动态流媒体内容的精准情感定位与分析。

最后，本研究主要聚焦于图文场景下的细粒度情感相关性解耦方法探索，并未对更广泛的下游任务进行实际应用。尽管已有的实验证明了多角度相关性建模对于提升情感分析性能的有效性，但是对于具体且复杂的多模态任务，如多模态反讽识别、跨模态虚假新闻检测以及面向开放域的舆情监控等，本研究的方法和模型还需要进一步验证。因此，未来的工作将致力于将本研究提出的模型框架迁移应用到这些具体的下游任务中，通过在真实且充满噪声的数据环境中进行测试，以全面验证其泛化能力与实用价值。